\documentclass[11pt]{article}
\usepackage[T1]{fontenc}
\usepackage{lmodern,amsmath,amssymb,amsthm,booktabs,microtype,xcolor}
\usepackage[letterpaper,margin=1in]{geometry}
\usepackage[colorlinks=true,linkcolor=blue!45!black,urlcolor=blue!45!black]{hyperref}
\newtheorem{theorem}{Theorem}[section]
\newtheorem{lemma}[theorem]{Lemma}
\newtheorem{proposition}[theorem]{Proposition}
\newtheorem{corollary}[theorem]{Corollary}
\newtheorem{remark}[theorem]{Remark}
\setlength{\emergencystretch}{3em}
\title{The eventual exact maximum number of cycles}
\author{Collaborative research working draft}
\date{11 November 2026}
\begin{document}
\maketitle
\begin{abstract}
Let $\psi(r)$ be the maximum number of cycles in a finite graph of cycle rank $r$.
Using the established unrestricted sharp-order proof and terminal-path estimates, we obtain, for all sufficiently large ranks,
\[
\begin{aligned}
 \psi(2q+1)&=2^{2q}+(2q+\tfrac12)2^q-3q,\\
 \psi(2q+2)&=2^{2q+1}+(2q+\tfrac92)2^q-3q-4.
\end{aligned}
\]
These are the counts attained by the Horak--Kim rings.
The new local step gives a fixed gap outside two saturated four-terminal profiles, uniformly over unbalanced cuts.
Exact ring optimization handles the remaining profile.
An exact running-maximum recurrence transfers the local bound to all graphs.
In particular, the normalized second-order coefficients are $1/\sqrt2$ on odd ranks and $1/2$ on even ranks.
\end{abstract}

\paragraph{Proof status.}
The new ring formulas, rigidity argument, and transfer recurrence have separate line audits by the root, four-regular, path-stability, and random-host agents.
Finite exact checks corroborate the formulas and case analysis; they do not replace the proofs.
The theorem has not been formalized in Lean or externally peer reviewed.
The absolute rank threshold is not evaluated.
The stable source modules used by Fable's formalization are unchanged.
The lower-bound construction and its cycle counts are due to Horak and Kim~\cite{HK}; the new claim is eventual unrestricted optimality.

\clearpage
\tableofcontents

\section{Conventions and established inputs}
The cycle rank of $G$ is $|E(G)|-|V(G)|+\operatorname{comp}(G)$.
Cycles are simple circuits; loops and parallel pairs count as cycles of lengths one and two.
Parallel edges are distinguished.
Subdivision preserves both cycle rank and cycle count, so allowing multigraphs does not change $\psi$.
The established cubic reduction gives a connected cubic multigraph attaining $\psi(r)$ at each relevant rank.
A cyclic cut has a cycle on each shore.

A normalized $k$-pole is connected, with $k$ distinct degree-two terminals and degree three elsewhere.
Its terminal paths are simple paths.
For a four-pole of rank $h$, divide all path counts $P_{ij}$ and disjoint path-pair counts $f_\alpha$ by $2^h$, writing $a_{ij}$ and $u_\alpha$.
Here $\alpha$ runs over the three pairings of the four terminals.
Put $S=\sum a_{ij}$ and $F=\sum u_\alpha$.
The affine cycle-space bounds give $0\le a_{ij}\le1$ and $0\le F\le1$.

This packet gives the new deductions in full, while taking the following previously audited results as inputs.
Their proof sources remain in the neighboring \nolinkurl{2026-11-10} directory.

\begin{theorem}[Established global inputs]
\label{thm:exact-packet-global-inputs}
There are absolute constants and thresholds with these properties.
\begin{enumerate}
\item $\psi(r)\le2^{r-1}+Kr2^{r/2}$.
\item At sufficiently large rank and sufficiently large fixed girth, every edge of a shortest cycle in a connected cubic graph belongs to at most $2^{r-2}$ cycles.
\item For each fixed girth, a connected cubic graph without a cyclic cut of size at most four has fewer than $2^{r-1}$ cycles at sufficiently large rank.
\item A connected loopless cubic graph with a cyclic cut of size at most three has at most $2^{r-1}+C2^{r/2}$ cycles.
\end{enumerate}
\end{theorem}

Items one through three are proved in \nolinkurl{unrestricted-sharp-cycle-bound.tex}; item four is in the small-cut proof modules assembled by \nolinkurl{small-cut-sharp-bound.tex}.
Only the exact shortest-cycle edge estimate and fixed-girth deficit are needed for the new running-maximum recurrence.

\begin{theorem}[Established path saturation gap]
\label{thm:terminal-path-saturation-gap}
Each normalized path coordinate of a loopless cubic two-pole is either one or at most $15/16$.
The conclusion applies to selected terminal pairs in three- and four-poles after suppressing the unused terminals.
\end{theorem}
Source: \nolinkurl{second-order-path-gap.tex}.

\begin{theorem}[Established high-path consequences]
\label{thm:three-pole-two-high}
For every fixed $\epsilon>0$, a normalized two-pole of rank $h$ with path density at least $1/2+\epsilon$ has at most $C_\epsilon2^{h/2}$ internal cycles.
In a three-pole, two path coordinates sharing a terminal and exceeding $1/2+\epsilon$ force the third coordinate to be at most $C'_\epsilon2^{-h/2}$.
The same statements hold for selected pairs in a four-pole after suppressing unused terminals.
\end{theorem}
Sources: the high-path cycle tradeoff and \nolinkurl{second-order-three-pole-tradeoff.tex}.

\begin{lemma}[Established one-full-path consequence]
\label{lem:one-full-tripole-small-adjacent}
In a rank-$h$ four-pole, $a_{12}=1$ implies
\[
 \min(a_{13},a_{23}),\ \min(a_{14},a_{24})\le3\,2^{-h/4}.
\]
\end{lemma}
Source: \nolinkurl{second-order-one-full-tripole.tex}.

\begin{corollary}[Established cross-sum forest estimate]
\label{cor:linear-four-pole-cross-sum-forest}
For each fixed $\eta>0$, a rank-$h$ four-pole with
$\sum_{i,j=1}^2a_{L_iR_j}\ge2+\eta$ satisfies
$F\le A_\eta(h+3)2^{-h/2}$.
\end{corollary}
Source: \nolinkurl{second-order-two-cycle-packing-refinement.tex}.

\begin{lemma}[Established added-edge and low-state bounds]
\label{lem:four-pole-added-edge}
For sufficiently large $h$, if both coordinates of an opposite pairing $\alpha=ij|kl$ exceed $0.01$, then
\[
 \max(a_{ij},a_{kl})+F-u_\alpha\le1.01.
\]
If all six path coordinates are at most $0.51$, then
\[
 S+F\le3.82,\qquad
 S+\tfrac12F+\tfrac12u_\gamma\le3.7
\]
for each selected pairing $\gamma$.
\end{lemma}
Source: \nolinkurl{second-order-all-four-cuts.tex}.

\begin{lemma}[Established saturated four-pole classification]
\label{lem:saturated-four-pole-double-cap}
Suppose all four path coordinates across a selected partition into left and right terminal pairs equal one.
If there is a pair of vertex-disjoint left-to-right paths, the pole is a chain of square and single-rung tiles, with the conventions in the next section.
If there is no such pair, the pole has an internal bridge separating its left terminals from its right terminals, and both resulting components have positive rank.
\end{lemma}
Source: \nolinkurl{second-order-saturated-blocks.tex}, including the zero-through case in the proof of its double-cap lemma.

Finally, in the absence of smaller cyclic cuts, a cyclic four-cut can be normalized to connected four-poles with distinct terminals and positive ranks $h,k$, satisfying $r=h+k+3$.
The exact crossing formula is stated and used in Section~\ref{sec:four-coefficient-rigidity}.
If $h=1$, the smaller shore is $C_4$; its normalized forest vector gives precisely the second low-state expression in the preceding lemma.
These normalizations are part of \nolinkurl{second-order-all-four-cuts.tex}.

\section{Exact optimization of square and single-rung rings}
\label{sec:four-coefficient-rings}

A square tile has two left vertices and two right vertices, with all four edges between the two pairs.
A single-rung tile has two vertices joined by one edge; these vertices serve as both its left and right pair.
Join the right pair of each tile to the next left pair by two matching edges, cyclically.
Each interface matching may be straight or crossed.
Let $a$ count square tiles and $s$ count single-rung tiles.
Put $W=2a+s$ and $m=a+s$.
The resulting cubic graph has rank $r=W+1$.
We restrict to $W\ge7$, so the small one-tile and two-tile degeneracies do not arise.

\begin{proposition}[Ordering the single-rung tiles]
\label{prop:four-coefficient-ring-order}
Suppose $a\ge1$.
Among all cyclic orders of these tiles, the maximum cycle excess over $2^W$ is
\[
 E(a,s)=2^a\left(2a+4s+\binom s2+\frac12\right)
        -3a-4s+\binom s2.
\]
It is attained when all single-rung tiles are consecutive.
The count is independent of the interface twists.
\end{proposition}

\begin{proof}
Assign weight $v_i=1/2$ to a square tile and $v_i=1$ to a single-rung tile.
For ordered tile positions $i,j$, let $v[i,j)$ be the product along the forward cyclic interval from $i$ through the position before $j$.
When $i=j$, use the empty interval and weight one.
Define
\[
 Q=\sum_{i,j=1}^{m}v[i,j).
\]

Cut parity makes every interface count odd, or makes all of them even.
The odd case uses one edge at each interface.
There are $2^m$ interface choices and two local paths in each square, hence exactly $2^{m+a}=2^W$ cycles.
A local path in a single-rung tile may have length zero when both interfaces choose the same vertex.

An even-sector cycle with no used interface is one of the $a$ squares.
If it uses every interface twice, each square has two through matchings.
Exactly half their $2^a$ choices make one circuit, giving $2^{a-1}$ cycles.
This remains true with any interface twists.

Otherwise the used interfaces form one nonempty proper cyclic interval.
Its two end tiles supply caps and its inner tiles supply through pairs.
A square has two choices in either role, and a single-rung tile has one.
The count is therefore $2^a$ times the product of the weights of the omitted tiles.
The omitted intervals have lengths zero through $m-2$.
The length-$m-1$ intervals, which are included in $Q$, have total weight
\[
 2^{-a}\sum_i v_i^{-1}=2^{-a}(2a+s)=2^{-a}W.
\]
Thus the cycle excess is exactly
\[
 a+2^{a-1}+2^aQ-W
 =2^a(Q+1/2)-m.
\]

Decompose $Q$ by the types of its two endpoints.
The square--square contribution is $2a(1-2^{-a})$.
For each single-rung tile, the two orders of mixed endpoints contribute
$2(1-2^{-a})$ and $1-2^{-a}$.
The equal single-rung endpoints contribute $s$.
For two distinct single-rung positions, their two cyclic interval weights are $z$ and $2^{-a}/z$, with $2^{-a}\le z\le1$.
Convexity gives
\[
 z+2^{-a}/z\le1+2^{-a}.
\]
Equality holds when one of those intervals contains no square.
All these equalities hold together when the single-rung tiles are consecutive.
Consequently
\[
 Q\le2a+4s+\binom s2
       +2^{-a}\left(-2a-3s+\binom s2\right).
\]
Substitution proves the formula and attainment.
\end{proof}

\begin{theorem}[The exact ring maximum at fixed rank]
\label{thm:four-coefficient-ring-maximum}
For $r\ge8$, every square and single-rung ring of rank $r$ has at most $H(r)$ cycles, where
\[
 \begin{aligned}
 H(2q+1)&=2^{2q}+(2q+1/2)2^q-3q,\\
 H(2q+2)&=2^{2q+1}+(2q+9/2)2^q-3q-4.
 \end{aligned}
\]
The maxima use no single-rung tile when $r$ is odd, and one when $r$ is even.
These are the Horak--Kim rings.
\end{theorem}

\begin{proof}
Fix $W=r-1\ge7$.
For $a\ge2$, compare $E(a,s)$ with $E(a-1,s+2)$, keeping $2a+s=W$.
Direct subtraction gives
\[
 E(a-1,s+2)-E(a,s)
 =-2^{a-1}\left(W+\frac{s^2+s}{2}-\frac{13}{2}\right)+2s-4<0.
\]
For $s=0,1$, the last inequality is immediate.
For $s\ge2$, the parenthesis exceeds $2s-4$ when $W\ge7$.
Thus among rings with at least one square, the maximum has $s=0$ or $s=1$ according to the parity of $W$.
Proposition~\ref{prop:four-coefficient-ring-order} gives the displayed formulas.

If $a=0$, every proper interval of used interfaces has one completion.
There are $W(W-1)$ such intervals.
The all-used case contributes at most one additional cycle.
Thus its excess is at most $W(W-1)+1$.
The stated maximum exceeds this quantity at $W=7,8$.
Its excess more than doubles when $W$ increases by two, whereas
\[
 (W+2)(W+1)+1<2\bigl(W(W-1)+1\bigr)\qquad(W\ge7).
\]
Induction on each parity proves the comparison for all $W\ge7$.
\end{proof}

\begin{remark}[A cap formula for open chains]
For an open chain with the same tile types, the through count is $D=2^a$.
Its disjoint left-cap/right-cap count is
\[
 C=2^a\sum_{i<j}2^{-a(i,j)},
\]
where $a(i,j)$ counts square tiles strictly between positions $i,j$.
The turn tiles contribute their local cap counts; all tiles before the first turn and after the second turn contribute through counts.
This gives the formula directly.
In particular,
\[
 C\le2^a\left(2a-4+2^{2-a}+4s+\binom s2\right).
\]
The square--square sum is exact, each single-rung tile contributes at most four against all squares, and each pair of single-rung tiles contributes at most one.
This sharpens the earlier bound to leading coefficient one when the number of single-rung tiles is minimal.
\end{remark}

\paragraph{Status.}
These formulas are new deductions from the previously audited saturated-ladder classification.
They have passed independent line audits by the root, path-stability, and random-host agents.
The companion checker verifies 527520 even states in 456 rings, 86 open chains, and 62491 exact optimization inequalities.
The proofs remain separate from the stable proof packet.

\section{Four-pole rigidity and the eventual small-cut maximum}
\label{sec:four-coefficient-rigidity}

This argument uses previously audited terminal inequalities.
It does not change their proofs or constants.
For a rank-$h$ four-pole, normalize all path and forest counts by $2^h$.
Write $a_{ij}$ for its six path counts and $u_\alpha$ for its three forest-pairing counts.
Put $S=\sum a_{ij}$ and $F=\sum u_\alpha$.
Each path coordinate and $F$ belongs to $[0,1]$.

We use the following established consequences, with absolute constants.
\begin{enumerate}
\item Every path coordinate is either one or at most $15/16$
      (Theorem~\ref{thm:terminal-path-saturation-gap}).
\item Two path coordinates sharing a terminal and exceeding $0.51$ force the third coordinate to be $O(2^{-h/2})$
      (Theorem~\ref{thm:three-pole-two-high}).
\item If $a_{12}=1$, then each of $\min(a_{13},a_{23})$ and $\min(a_{14},a_{24})$ is at most $3\,2^{-h/4}$
      (Lemma~\ref{lem:one-full-tripole-small-adjacent}).
\item A cross sum at least $2.01$ gives $F=O((h+1)2^{-h/2})$
      (Corollary~\ref{cor:linear-four-pole-cross-sum-forest}).
\item If both coordinates of an opposite pair $\alpha$ exceed $0.01$, then
      $\max(a_{ij},a_{kl})+F-u_\alpha\le1.01$ for sufficiently large $h$
      (Lemma~\ref{lem:four-pole-added-edge}).
\end{enumerate}
Suppressing unused ports gives these consequences for four-poles.
A loop produced by suppression would force the relevant path densities to be at most one half, so it does not affect the high-density cases.

\begin{lemma}[A fixed gap outside two saturated profiles]
\label{lem:four-coefficient-aggregate-rigidity}
There is an absolute $h_0$ such that every loopless normalized four-pole of rank $h\ge h_0$ with
\[
 S+F\ge3.99
\]
has either a full three-edge star or a full $K_{2,2}$ among its path coordinates.
Here an edge is full when its path count equals $2^h$.
\end{lemma}

\begin{proof}
Increase $h_0$ until every two-high small coordinate and every one-full minimum in the preceding list is at most $0.01$.
Also require the forest bound in item four to be at most $0.01$.

First suppose some cross sum is at least $2.01$.
Then $F\le0.01$, so $S\ge3.98$.
Form a graph on the four ports using edges with $a_{ij}>\tau=0.51$.
It is triangle-free, since two such adjacent edges make the third coordinate at most $0.01$.

If this high graph is $K_{2,2}$ but one edge is not full, then
\[
 S\le3+15/16+2(0.01)=3.9575<3.98.
\]
If it has three edges, it is a star or a path.
Their sums are at most $3.03$ and $3+0.51+0.02=3.53$, respectively.
Two adjacent high edges give $S\le2+0.01+3(0.51)=3.54$.
At most one high edge gives $S\le1+5(0.51)=3.55$.

The remaining case consists of two disjoint high edges.
If neither is full, then
\[
 S\le2(15/16)+4(0.51)=3.915.
\]
If one is full, item three makes two distinct cross coordinates at most $0.01$.
The other two cross coordinates are at most $0.51$.
Thus $S\le2+2(0.51)+2(0.01)=3.04$.
All alternatives contradict $S\ge3.98$.
The four high edges must therefore form a full $K_{2,2}$.

Now suppose every cross sum is less than $2.01$.
The three cross sums add to $2S$, so $S<3.015$.
The hypothesis gives $S\ge2.99$ and $F>0.975$.
Every opposite pair of coordinates has sum greater than $2.99-2.01=0.98$.

If no coordinate is full, each is at most $15/16$ and consequently greater than $0.98-15/16=0.0425$.
Apply item five to all three opposite pairings.
Their sum gives $S/2+2F\le3.03$, hence $S+4F\le6.06$.
But $S\ge2.99$ and $F>0.975$ give $S+4F>6.89$, a contradiction.

There is therefore a full coordinate, say $a_{12}=1$.
Item three gives one small coordinate in each of $\{a_{13},a_{23}\}$ and $\{a_{14},a_{24}\}$.
Each opposite coordinate exceeds $0.98-0.01=0.97$, so the saturation gap makes it full.
The two small coordinates cannot be opposite, because their opposite-pair sum exceeds $0.98$.
The three full coordinates therefore form a star.
\end{proof}

\begin{lemma}[Aligned saturated cross profiles]
\label{lem:four-coefficient-aligned-rings}
Glue two saturated four-poles across their four ports, with their full $K_{2,2}$ partitions aligned.
If either pole has no through pair, the resulting cubic graph has a cyclic cut of size three.
Otherwise the result is a square and single-rung ring.
\end{lemma}

\begin{proof}
With positive through count, the saturated four-pole classification gives a chain of square and single-rung tiles.
The two aligned chains close into the ring described in Section~\ref{sec:four-coefficient-rings}.

For completeness, the zero-through alternative in the proof of Lemma~\ref{lem:saturated-four-pole-double-cap} gives a bridge inside that pole.
It separates its left terminal pair from its right terminal pair, with positive rank on each side.
One bridge shore has two external port edges in the glued graph.
Its boundary thus consists of those two edges and the bridge.
Both global shores are connected and cyclic.
This is a cyclic three-edge cut.
\end{proof}

\begin{theorem}[Eventual exact maximum in the small-cut class]
\label{thm:four-coefficient-eventual-small-cut}
There is an absolute rank threshold $R$ such that every connected loopless cubic multigraph of rank $r\ge R$ with a cyclic cut of size at most four satisfies
\[
 c(G)\le H(r),
\]
where $H$ is the Horak--Kim count in Theorem~\ref{thm:four-coefficient-ring-maximum}.
The bound is attained for every $r\ge R$.
\end{theorem}

\begin{proof}
The established theorem for cyclic cuts of size at most three has error $O(2^{r/2})$.
That error is eventually smaller than $H(r)-2^{r-1}$.
We may therefore assume no such cut exists.

Normalize a cyclic four-edge cut as in the previously audited four-cut proof.
Its shores are connected four-poles with distinct ports and positive ranks $h\le k$.
They satisfy $r=h+k+3$.
Call the poles $X,Y$, with $Y$ of rank $k$.
In particular $k\to\infty$ uniformly as $r\to\infty$.
The exact normalized crossing count is
\[
 \mathcal B(X,Y)
 =\sum_{i<j}a_{ij}(X)a_{ij}(Y)
  +\sum_\alpha u_\alpha(Y)\bigl(F(X)-u_\alpha(X)\bigr).
\]
Every coefficient from $X$ is at most one, so $\mathcal B(X,Y)\le S(Y)+F(Y)$.

First suppose all six path coordinates of $Y$ are at most $0.51$.
The existing low-state lemma gives $S(Y)+F(Y)\le3.82$.
The asymptotic cycle bound gives $c(Y)\le(9/16)2^k$ for large $k$.
If $h\ge2$, then
\[
 \frac{c(G)}{2^{h+k}}
 \le3.82+\frac9{16}2^{-h}+2^{-k}
 \le3.960625+2^{-k}<4.
\]
If $h=1$, then $X=C_4$.
Its exact forest vector instead gives the low-state bound $\mathcal B(X,Y)\le3.7$.
Therefore
\[
 \frac{c(G)}{2^{h+k}}
 \le3.7+\frac9{32}+2^{-k}<4
\]
for sufficiently large $k$.
These bounds are uniform in the smaller rank.

Now suppose a path coordinate of $Y$ exceeds $0.51$.
The high-path cycle tradeoff gives $c(Y)\le C2^{k/2}$ for an absolute $C$.
If $S(Y)+F(Y)<3.99$, then
\[
 \frac{c(G)}{2^{h+k}}
 \le3.99+C2^{-h-k/2}+2^{-k}<4
\]
for large $k$, uniformly for every $h\ge1$.
We may assume $S(Y)+F(Y)\ge3.99$ and apply the rigidity lemma.

Suppose first that $Y$ has a full star.
Its other three path coordinates are $O(2^{-k/2})$.
If $X$ misses a full coordinate on this star, the saturation gap loses at least $1/16$ in the crossing sum.
Since $F(Y)\le1$, this gives
\[
 \mathcal B(X,Y)\le4-1/16+O(2^{-k/2}).
\]
The same normalized internal-cycle bounds make $c(G)<2^{r-1}$ for large $k$.

If $X$ is full on the same star, its three leaf-pair coordinates are $O(2^{-h/2})$.
The three full path products contribute exactly three.
The forest term is at most $F(X)F(Y)\le1$.
The remaining path products contribute $O(2^{-(h+k)/2})$.
Both poles have full paths, so their internal cycle counts are $O(2^{h/2})$ and $O(2^{k/2})$.
Consequently
\[
 c(G)\le2^{r-1}+C'2^{r/2}
\]
with an absolute $C'$, including all bounded smaller ranks.
This is eventually less than $H(r)$.

Finally suppose $Y$ has a full $K_{2,2}$.
Its remaining path coordinates are $O(2^{-k/2})$ and its forest total is $O((k+1)2^{-k/2})$.
If $X$ misses any of these four full coordinates, then
\[
 \mathcal B(X,Y)\le4-1/16+O((k+1)2^{-k/2}).
\]
The internal counts again give $c(G)<2^{r-1}$ for large $k$, uniformly in $h$.
Otherwise $X$ and $Y$ are aligned saturated cross poles.
The zero-through case in Lemma~\ref{lem:four-coefficient-aligned-rings} contradicts the absence of a cyclic three-edge cut.
The other case is a square and single-rung ring, so Theorem~\ref{thm:four-coefficient-ring-maximum} gives $c(G)\le H(r)$.

All thresholds and implied constants used above are absolute.
Their maximum supplies $R$.
The Horak--Kim rings have cyclic four-edge cuts and attain $H(r)$ on both parities.
\end{proof}

\paragraph{Status.}
This refinement has passed independent line audits by the root, path-stability, and random-host agents.
Its proof treats the smaller shore without hiding its rank in a constant.
The companion checker tests 1059 exact rational endpoint profiles; their maximum path sum outside a full $K_{2,2}$ is $1583/400=3.9575$.
The eventual rank threshold is not made effective here.

\section{Exact and parity-resolved transfer from cyclic four-cuts}
\label{sec:new-coefficient-transfer}

These reductions use the completed unrestricted sharp-order proof.
They do not require a better estimate for its unspecified constant.
All maxima below range over finite graphs, with parallel edges distinguished.

Put
\[
 E_r=\psi(r)-2^{r-1}.
\]
Let $B_r$ be the maximum excess among connected loopless cubic multigraphs of rank $r$ with a cyclic cut of size at most four.
Set $B_r=0$ if this class is empty.

\begin{proposition}[An exact running-maximum reduction]
\label{prop:new-excess-recurrence}
For all sufficiently large $r$,
\[
 E_r\le\max\{B_r,E_{r-1},0\}.
\]
Consequently, for one fixed integer $R$ and every $r\ge R$,
\[
 E_r\le\max\{E_{R-1},0,B_R,B_{R+1},\ldots,B_r\}.
\]
\end{proposition}

\begin{proof}
Choose a cubic multigraph attaining $\psi(r)$, as in the established reduction.
If it has a loop, deleting that loop gives
\[
 E_r\le E_{r-1}-2^{r-2}+1\le E_{r-1}.
\]
If it has a cyclic cut of size at most four, use the definition of $B_r$.

Otherwise, the bounded-girth deficit excludes all girths below the fixed threshold in the exact-half forest theorem, at sufficiently large rank.
Such graphs have negative excess.
In the remaining high-girth case, an edge of a shortest cycle has at most $2^{r-2}$ incident cycles.
Deleting that edge leaves a connected rank-$(r-1)$ graph, so
\[
 \psi(r)\le\psi(r-1)+2^{r-2},
 \qquad E_r\le E_{r-1}.
\]
All thresholds are absolute and independent of $B_r$.
Iteration proves the second assertion.
\end{proof}

\begin{corollary}[Transfer of an eventual exact formula]
\label{cor:new-eventual-exact-transfer}
Suppose graphs of rank $r$ attain $2^{r-1}+H_r$ for every sufficiently large $r$.
Assume that $H_r$ is eventually nondecreasing and tends to infinity.
If $B_r\le H_r$ for all sufficiently large $r$, then
\[
 \psi(r)=2^{r-1}+H_r
\]
for all sufficiently large $r$.
\end{corollary}

\begin{proof}
Increase $R$ so that all hypotheses and Proposition~\ref{prop:new-excess-recurrence} hold above $R$.
The running-maximum bound gives
\[
 E_r\le\max\{E_{R-1},0,H_r\}.
\]
Eventually $H_r$ exceeds the two fixed initial terms.
The assumed construction gives the reverse inequality.
\end{proof}

For the Horak--Kim construction, the excess is
\[
 H_r=
 \begin{cases}
 (r-\tfrac12)2^{(r-1)/2}-\tfrac32(r-1),&r\text{ odd},\\[1mm]
 (r+\tfrac52)2^{(r-2)/2}-\tfrac12(3r+2),&r\text{ even}.
 \end{cases}
\]
The two alternating increments are
\[
 H_{2q+2}-H_{2q+1}=4\,2^q-4,
 \qquad
 H_{2q+3}-H_{2q+2}=(2q+\tfrac12)2^q+1.
\]
They are positive for $q\ge1$, and $H_r\to\infty$.
Thus an eventual sharp upper bound for the small-cut class would transfer as an exact formula, not merely a leading coefficient.

\begin{proposition}[The two parity limits]
\label{prop:new-parity-limsup}
Normalize the unrestricted and small-cut excesses by
\[
 a_r=\frac{E_r}{r2^{r/2}},\qquad
 b_r=\frac{\max\{B_r,0\}}{r2^{r/2}}.
\]
Let $A_o,A_e$ and $B_o,B_e$ be their respective odd and even limsups.
Then
\[
 B_o\le A_o\le\max\{B_o,B_e/\sqrt2\},
 \qquad
 B_e\le A_e\le\max\{B_e,B_o/\sqrt2\}.
\]
In particular, $\limsup a_r=\limsup b_r$.
\end{proposition}

\begin{proof}
The established sharp-order bound makes all four limsups finite.
The known constructions make the unrestricted excess positive eventually.
The recurrence gives
\[
 A_o\le\max\{B_o,A_e/\sqrt2\},
 \qquad
 A_e\le\max\{B_e,A_o/\sqrt2\}.
\]
Substitution yields
\[
 A_o\le\max\{B_o,B_e/\sqrt2,A_o/2\}.
\]
Since $A_o>0$, it follows that $A_o\le\max\{B_o,B_e/\sqrt2\}$.
The even bound is the same.

The direct inequalities $A_o\ge B_o$ and $A_e\ge B_e$ follow from inclusion of the small-cut class.
Taking the maximum of the two upper bounds gives
$\max\{A_o,A_e\}\le\max\{B_o,B_e\}$.
The reverse inequality follows from inclusion, proving the global limsup equality.
\end{proof}

The inequalities do not assert equality separately for each parity.
A crossed lower transfer would need an additional graph construction.
The proved bounds already transfer the proposed parity coefficients whenever the local upper coefficients are $B_o\le1/\sqrt2$ and $B_e\le1/2$.
The Horak--Kim lower bounds then give the matching unrestricted limits.

\section{The eventual exact unrestricted maximum}
\label{sec:eventual-exact-cycle-theorem}

\begin{theorem}[Eventual exact Horak--Kim formula]
\label{thm:eventual-exact-cycle-maximum}
There is an absolute integer $R$ such that the maximum number $\psi(r)$ of cycles in a finite graph of cycle rank $r\ge R$ is
\[
\begin{aligned}
 \psi(2q+1)&=2^{2q}+(2q+\tfrac12)2^q-3q,\\
 \psi(2q+2)&=2^{2q+1}+(2q+\tfrac92)2^q-3q-4.
\end{aligned}
\]
The Horak--Kim rings attain both formulas.
\end{theorem}

\begin{proof}
Let $H(r)$ denote the two displayed counts and put $H_r=H(r)-2^{r-1}$.
Theorem~\ref{thm:four-coefficient-eventual-small-cut} bounds the maximum small-cut excess $B_r$ by $H_r$ at every sufficiently large rank.
The exact ring count supplies a graph attaining $H(r)$ at each such rank.
The two increments computed in Section~\ref{sec:new-coefficient-transfer} show that $H_r$ is eventually strictly increasing and tends to infinity.
Corollary~\ref{cor:new-eventual-exact-transfer} therefore gives the claimed unrestricted equality.
\end{proof}

\begin{corollary}[Sharp second-order coefficients]
\[
 \psi(r)=2^{r-1}+(1+o(1))r2^{\lfloor(r-1)/2\rfloor}.
\]
More precisely,
\[
 \lim_{\substack{r\to\infty\\r\text{ odd}}}
 \frac{\psi(r)-2^{r-1}}{r2^{r/2}}=\frac1{\sqrt2},
 \qquad
 \lim_{\substack{r\to\infty\\r\text{ even}}}
 \frac{\psi(r)-2^{r-1}}{r2^{r/2}}=\frac12.
\]
\end{corollary}

\begin{proof}
Divide the exact formulas by the indicated scales.
The additive terms linear in $r$ vanish after normalization.
\end{proof}

\paragraph{Scope.}
The threshold is not evaluated.
The result does not assert the Horak--Kim formula at every small rank, and does not settle the relative edge-half conjecture.
The proof is a new mathematical working proof, with separate internal audits, based on the established source modules listed in this packet.
It has not been formalized in Lean or externally peer reviewed.

\section{Eventual uniqueness among cubic maximizers}
\label{sec:four-coefficient-uniqueness}

\begin{corollary}
\label{cor:four-eventual-cubic-uniqueness}
There is an absolute integer $R$ such that, for every $r\ge R$, a connected cubic multigraph $G$ of rank $r$ satisfies
\[
 c(G)=\psi(r)
\]
if and only if $G$ is isomorphic to the Horak--Kim graph $MC_{r-1}$.
For odd $r$, this is the ring of $(r-1)/2$ square tiles.
For even $r$, it is the ring of $(r-2)/2$ square tiles and one single-rung tile.
The assertion concerns cubic graphs; it does not classify arbitrary subdivisions or other noncubic maximizers.
\end{corollary}

\begin{proof}
Increase $R$ until the eventual exact formula and every threshold in its proof apply.
Write $H_r=\psi(r)-2^{r-1}$ for the Horak--Kim excess.
Its two alternating increments are strictly positive.

The loop and high-girth deletion branches give excess at most $H_{r-1}<H_r$.
The bounded-girth branch has negative excess.
Thus every cubic maximizer has a cyclic cut of size at most four.
Cyclic cuts of size at most three give only $O(2^{r/2})$ excess, which is strictly smaller than $H_r$ for large $r$.

In the four-cut proof, every fixed-deficit branch lies below one half.
The aligned full-star branch again has only $O(2^{r/2})$ excess.
The only remaining case is an aligned pair of saturated cross poles with positive through counts.
It is a square and single-rung ring.

The fixed-rank optimization in Theorem~\ref{thm:four-coefficient-ring-maximum} is strict whenever the number of single-rung tiles exceeds its parity minimum.
The all-single-rung case is also strictly smaller.
Thus a maximizing ring has no single-rung tile at odd rank and exactly one at even rank.
Its cyclic tile order is unique up to rotation.

Finally, interface twists do not change its isomorphism type.
Swapping both strand labels in a tile transfers a twist from one adjacent interface to the other.
A square tile additionally permits swapping its left pair while fixing its right pair, since its internal graph is $K_{2,2}$.
Move all twists to an interface next to a square and remove the remaining twist by this last automorphism.
There is at least one square at the ranks under consideration.
The untwisted ring is exactly $MC_{r-1}$.
Conversely, that graph attains the eventual exact formula.
\end{proof}

The rank threshold is unspecified, as in the eventual exact-count theorem.
No uniqueness assertion is made at smaller ranks.


\section{Exact checks and review scope}
The companion \nolinkurl{four-coefficient-check.py} uses exact integer and rational arithmetic.
Its saved results cover 527520 even affine states in 456 mixed rings, including both interface twists and all tile words through weight ten; 86 open-chain formulas; 62491 fixed-rank optimization inequalities; and 1059 rational endpoint profiles in the rigidity cases.
The largest path sum outside a full $K_{2,2}$ in the checked endpoint cases is $1583/400=3.9575$.

The line audits check the omitted-interval indexing, the all-through parity count, fixed-rank comparisons, the zero-through bridge alternative, constants uniform in the smaller shore, and the exact recurrence through rank deletion.
The audit record distinguishes the proved parity limsup inequalities from a stronger, unproved crossed lower-transfer assertion.
The eventual exact theorem uses the monotone-excess corollary and does not require such an assertion.

\begin{thebibliography}{9}
\bibitem{HK} P. Horak and J. Kim, \emph{On the maximum number of cycles in a graph}, Journal of Graph Theory 99 (2022), 671--690.
The explicit ring counts are in Theorem 5.1 of the revised preprint, \url{https://arxiv.org/html/2006.13459v2}.
\end{thebibliography}
\end{document}
